\section{What a formal legal language actually does}
\label{sec:language-lesson}

Consider the instruction to obtain written agreement before applying the
streamlined approach. A checklist can remind an officer to obtain it. A database
can store a scanned agreement. Neither tells a program what to conclude if the
agreement covers bonds while the proposed transaction is in a different category,
or if the client withdrew yesterday. A formal language supplies a precise way
to state the relevant conditions and the changes that events produce. Its
purpose here is to preserve an agreed interpretation closely enough that people
can inspect it and computers can execute and check it.

The word ``formal'' does not mean that the language resolves legal judgment.
It means that a well-formed statement has a specified computational meaning.
We can then distinguish disagreement about that meaning from a programming
mistake in carrying it out. The original circular, the adopted interpretation
and the formal statement must remain connected.

\subsection{From a sentence to a calculation with a defined meaning}

A programming language needs three ingredients. Its \emph{syntax} says which
statements are allowed. Its \emph{types} say which kinds of value can be combined.
Its \emph{semantics} says what each allowed statement does. For example,
\texttt{portfolio >= HKD(40000000)} has comparison syntax, compares two money
values, and means an inclusive monetary comparison. Comparing that portfolio
with a consent date is a type error. Using \texttt{>} instead of \texttt{>=} is
well-typed but changes the meaning at the boundary. A type checker catches the
first error; an independently chosen boundary example catches the second.

For the SPI financial condition, a readable formal statement is:
\begin{lstlisting}[caption={Illustrative controlled syntax for Annex 1 paragraph 3.1. This is a teaching notation, not Catala or Stipula source.}]
financial_condition(client, date) :=
    portfolio(client, date) >= HKD(40000000)
    OR net_assets_ex_home(client, date) >= HKD(80000000)
\end{lstlisting}
Here \texttt{client} and \texttt{date} are arguments: every occurrence refers to
the same person and valuation date. \texttt{portfolio} and
\texttt{net\_assets\_ex\_home} name defined inputs. The source, paragraphs
3.2--3.4, determines the ownership and net-asset definitions; the host bank must
provide values calculated under the adopted mapping. The notation makes no
assumption that the portfolio is simply assets held at this bank.

The compiler stores the statement as a tree. The root is \texttt{OR}; its two
children are inclusive comparisons; each comparison has an input and a monetary
constant. This tree is called an \emph{abstract syntax tree}: it retains the
operations and their arguments without depending on the layout of the original
sentence. The JSON in Section~\ref{sec:product} is a serialized version of such
a tree. An intermediate representation, or IR, adds stable identifiers, source
locations, types and interpretation records to that executable structure.

Evaluation works from the leaves towards the root. For a portfolio of HK\$35
million and net assets of HK\$90 million, the comparisons give false and true;
their alternative gives true. Missing portfolio evidence produces an unresolved
comparison. An established net-asset route can still make the alternative true.
If neither route is established and one remains unresolved, the answer remains
unresolved. This is why an ordinary Java \texttt{boolean}, which has only two
values, is insufficient as the complete result type.

The financial predicate answers one question. It does not produce a duty or an
authorization to execute an order. The larger model must combine qualification,
the selected product category, consent, exposure and retained requirements.
Keeping these named subquestions separate is useful both to the reviewer and
to the engineer diagnosing a rejected request.

\subsection{Catala: general rules and explicit exceptions}

Catala was designed around a recurring structure in legislation: a value is
defined by a general rule, subject to exceptions. Its central contribution is
to give that structure an explicit executable meaning, while placing code close
to the legal text \citep[secs.~2--4]{catala2021}. It is useful when nested
qualifications would otherwise become a long, undocumented sequence of
\texttt{if} statements.

For a small banking illustration, take the product-explanation relief in Annex
1 paragraph 10.3. Within the selected solicited-transaction profile, an
explanation remains required when the client requests it or raises material
queries. A local representation can say:
\begin{lstlisting}[caption={An exception with an explicit guard. The combined guard avoids creating two competing exceptions for the same requirement.}]
normally: product_explanation_required = false
exception when client_requests_explanation
               OR client_has_material_queries:
    product_explanation_required = true
always retain: current_offering_documents_required = true
\end{lstlisting}
The \emph{guard} is the condition after ``when''. If the guard is true, use the
exception; if false, use the general value. If evidence about the guard is
missing, LegalMath leaves the requirement unresolved. The separate offering-
document requirement is not inside the explanation exception. A programmer
who implements one switch for both has changed the model
\citep[Annex 1, para.~10.3]{sfcspiannex1}.

There is an important distinction between borrowing this organization and
claiming to implement Catala itself. Catala's core calculus distinguishes an
empty applicable value and a conflict between exceptions. LegalMath additionally
needs incomplete bank evidence to remain explicitly unknown. Its proposed
default operator evaluates every sibling guard: two established exceptions
conflict; one established exception with an unresolved competing guard remains
unresolved; one established exception with all competitors refuted supplies the
value; all refuted guards use the general value. Nesting expresses an adopted
priority. File order does not.

Suppose two exceptions have been written, one for a request and another for a
material query. A client who does both would activate both. Under the proposed
sibling-conflict semantics, identical resulting values still conflict. Combining
the two conditions into one guard expresses the intended single exception.
That example shows why choosing a language is a substantive design decision:
a rule engine's treatment of overlapping rules must match the representation
the reviewer has adopted. The delivered Java example uses the equivalent
explicit Boolean condition for this one requirement; its compiler rejects the
general default operator until that operator has its own implementation tests.

\subsection{Stipula: agreements, authorized actions and state changes}
\label{sec:stipula-tutorial}

A wealth calculation is a function of facts. An agreement also has a history.
Before signature, after signature, after withdrawal and after an amendment, the
same attempted action can have different consequences. Stipula is a language
for expressing those interactions. The requested 2021 paper defines parties,
agreement terms, fields, assets, states, guarded functions and scheduled events
\citep[secs.~3--4]{stipula2021}. The original language in that paper is
type-free; the typed decision representation proposed here is a separate design.

A \emph{party} is an actor allowed to participate, such as a lender or borrower
in the paper's rental example. An \emph{agreement} establishes the parties and
the terms on which they agree. A \emph{field} stores information, such as a price
or agreed period. An \emph{asset} represents a transferable resource. A
\emph{state} records which stage the agreement has reached. A \emph{function}
specifies an action a named party may invoke in a particular state; a guard
adds a condition for the action. Its body changes fields, sends information or
transfers assets, and its transition identifies the next state. A scheduled
\emph{event} specifies an action to attempt at a future time if the agreement
is then in the required state.

For example, the rental contract in the paper first records an offer. The
borrower accepts by supplying the required payment, which moves the contract
into its using state. The contract provides the use information and schedules
an event for the end of the rental period. Either an appropriate borrower action
or the scheduled event can finish the interaction and release the held payment.
The language gives a precise account of which party can trigger each step and
when the payment moves. Assets have special transfer behavior; they are not
ordinary numeric fields that an assignment may freely duplicate
\citep[sec.~3, rental example and grammar]{stipula2021}.

The banking analogue uses the state and actor ideas but does not transfer client
assets. The following is explanatory pseudocode, deliberately distinguished
from compilable Stipula syntax:
\begin{lstlisting}[caption={A consent model inspired by Stipula. ``Record'' observes an event; ``allow'' describes the control's decision.}]
parties: Client, Bank
terms: client_id, product_category, threshold, monitoring_mode
states: NOT_ESTABLISHED, ACTIVE, WITHDRAWN

NOT_ESTABLISHED:
    record Client.written_agreement(terms, evidence)
    when Bank.confirmed_required_explanations(evidence)
    -> ACTIVE

ACTIVE:
    record Client.withdrawal(instruction_id) -> WITHDRAWN
    allow Bank.apply_streamlining(transaction)
    when category_matches(transaction, terms)
         AND all_other_required_conditions_hold(transaction)

WITHDRAWN:
    do not allow Bank.apply_streamlining(transaction)
    record a new, fully evidenced agreement -> ACTIVE
\end{lstlisting}
The terms bind consent to a client, category and arrangement. A new agreement
after withdrawal is a new generation of consent, not deletion of the withdrawal.
The regulatory duties themselves come from the applicable regulatory sources.
Client and bank agreement can establish the consent contemplated by those
sources; it cannot amend an SFC requirement. This distinction is essential when
borrowing a contract language for a regulatory-control system.
The figure makes those changes visible without requiring the reader to know a
particular language's punctuation.

\begin{figure}[htbp]
\centering
\begin{tikzpicture}[node distance=17mm, every node/.style={font=\small},
box/.style={draw=teal,fill=pale,rounded corners,align=center,text width=33mm,minimum height=14mm},
arr/.style={-{Latex[length=2mm]},thick,draw=navy}]
\node[box] (none) {Consent not established};
\node[box,right=of none] (active) {Active consent\\for these terms};
\node[box,right=of active] (withdrawn) {Withdrawn};
\draw[arr] (none)--node[above=5pt,align=center,font=\footnotesize]{Evidenced\\agreement}(active);
\draw[arr] (active)--node[above=5pt,font=\footnotesize]{Withdrawal}(withdrawn);
\draw[arr] (withdrawn.south) to[bend left=35] node[below]{New evidenced agreement} (active.south);
\end{tikzpicture}
\caption{Consent transitions for the example. An incomplete event history produces
an unresolved assessment outside these established states. It is not evidence
of an active agreement.}
\label{fig:consent-tutorial}
\end{figure}

Execution requires a state, stored terms and a history of ordered events. In
Stipula's operational semantics, the runtime also maintains pending scheduled
events and a clock. In the inspected 2021 semantics, due events take priority
over function calls at the same time. That is a language rule, not a rule of
Hong Kong banking law. LegalMath must adopt an explicit ordering policy for
withdrawal and order events, record the authority for their sequence, and retain
an unresolved result where the ordering cannot be established. It cannot import
Stipula's clock convention as if the circular prescribed it.

There is a second adaptation. A control may prohibit streamlined treatment after
withdrawal, but the bank might nevertheless have processed a transaction.
Deleting that transaction from the history because it is an illegal transition
would hide the very incident a compliance system needs to detect. LegalMath
therefore separates the observed event record from the decision about whether
the proposed action is permitted by the selected control. Normative languages
such as eFLINT and Symboleo help make this distinction explicit: a duty or
violation is represented separately from the fact that an action happened
\citep{eflint2020,symboleo2020}.

\subsection{What it means to verify such a model}

A property is a precise question about the model. For the consent example:
``After a known withdrawal, and before a new qualifying agreement, this control
never returns that the streamlining conditions are met.'' Testing checks selected
histories. Model checking explores all histories within a specified finite
model. A mathematical proof establishes the property under explicit assumptions.
None establishes that a real client signed a document, or that the selected
provisions are the complete legal regime.

Stipula's behavioral-equivalence work compares the observable interactions of
two formal contracts. Two implementations can use different internal state names
yet be equivalent if the relevant party actions, communications and asset
behavior match. This is useful for comparing formal implementations. It is not
a test that the English circular and the contract mean the same thing
\citep[sec.~5]{stipula2021}.

The Stipula--KeY work supplies a particularly relevant connection to Java.
Its translator builds Java verification models from a supported class of
Stipula contracts and adds specifications in the Java Modeling Language (JML).
KeY reasons about Java programs and these specifications. A JML precondition
states what must hold before a method; a postcondition states what the method
must establish; an invariant states a condition preserved at the designated
program points. For instance, an illustrative postcondition can require a
withdrawal method to leave consent inactive, while an invariant can prevent a
successful streamlining decision in that state. The proof obligation concerns
all executions covered by the model and its assumptions
\citep[secs.~3--5]{stipulakey2025,stipulakeycode}.

That translation is a useful research route, but its Java verification model is
not automatically a deployable bank integration. The paper restricts cycle
structure and abstracts some event conditions; the adapter still has to preserve
timestamps, evidence and the bank's actual interface. The concrete Java
demonstration in this proposal instead compiles the stated decision subset
directly. It compares Java with a separately implemented reference evaluator
and source-derived examples. A later KeY or Lean task can strengthen the
compiler evidence for a chosen property without changing who approves the
legal interpretation.

\subsection{Solvers, explanations and search in ordinary programming terms}

A satisfiability solver searches for values that make a set of logical
constraints true. To compare the correct wealth alternative with a mistaken
conjunction, give it both formulas and ask for inputs where their answers differ.
HK\$35 million and HK\$90 million are one such counterexample. An SMT solver
extends this idea with theories such as exact integer arithmetic. It is useful
here because money can be encoded as integers. Before treating ``no
counterexample'' as evidence, the system must establish that the allowed input
domain is nonempty and that both formulas were encoded faithfully. A timeout
means that the search did not finish.

A justification engine answers a related question: which rules and facts support
a conclusion? In a rule program, \texttt{eligible(C) :- financial(C), consent(C)}
means that both supporting predicates are needed to derive eligibility for the
same client \texttt{C}. s(CASP), used by s(LAW), can expose the supporting
derivation and assumptions rather than returning only a label. An assumed
consent predicate in a hypothetical answer is still an assumption; it must never
be written back as an observed agreement \citep{slaw2024}.

Search over interpretations sits above these calculations. An agent may propose
two readings of a clause, retrieve a referenced definition, and ask a solver for
a case that separates them. A Monte Carlo tree or UCB scheduler can choose which
investigation to try next. Its score measures usefulness under the chosen search
objective; it does not certify the resulting legal reading. The adopted meaning
must come from source-based adjudication. The existing MCP projects can assist
these bounded jobs, as explained in Section~\ref{sec:prototype}.
